%!TEX TS-program = xelatex
%!TEX encoding = UTF-8 Unicode
% Deedy CV/Resume Template
% Original author: Debarghya Das (http://debarghyadas.com)
% Repository: https://github.com/deedy/Deedy-Resume
% License: Apache 2.0
%
% IMPORTANT: Compile with XeLaTeX (requires fontspec + Lato/Raleway fonts)
%
% GENERATED FILE -- do not edit directly.
% Edit ../cv-data.yaml and this file's template (../templates/deedy-resume.tex.j2),
% then run `python3 generate_cvs.py` from the repo root.
%
% ============================================================
% NOTE(hlim): PHILOSOPHY OF THIS TEMPLATE
%
% Industry CVs (especially for big tech / AV companies) should be
% PROJECT-DRIVEN, not publication-driven.
%
% Recruiters ask: "What did you build, and did it actually work?"
% They do NOT read paper titles. Venue names (ICRA, NeurIPS) serve
% only as a quality signal — one line is enough for publications.
%
% The structure of this template reflects this:
%   - LEFT  column: Skills + open-source portfolio + Education
%   - RIGHT column: Experience (impact-driven) + Key Relevant Projects
%                   + Publications (1 line) + Awards
%
% For each job application, the "Key Relevant Projects" section (right col)
% should be CUSTOMIZED to match the job description.
% e.g., applying to an AV perception team → emphasize detection/segmentation.
%       applying to a mapping team        → emphasize SLAM/HD map projects.
%
% Reference: https://www.sweresume.app/articles/faang-resume-guide/
% ============================================================

\documentclass[]{deedy-resume-openfont}
\usepackage{fancyhdr}
\usepackage{fontawesome5}
\usepackage{pifont}
\newcommand{\githubstars}[1]{[{\faGithubSquare}\,\ding{72}\,\textbf{#1}]}

\pagestyle{fancy}
\fancyhf{}

\begin{document}

\lastupdated

%%%%%%%%%%%%%%%%%%%%%%%%%%%%%%%%%%%%%%
%
%     TITLE NAME
%
%%%%%%%%%%%%%%%%%%%%%%%%%%%%%%%%%%%%%%
% NOTE(hlim): Put your name, personal website, email, and one or two key links.
% Keep the contact line short — 3 items max. Google Scholar is useful for
% researchers; LinkedIn is useful for everyone. Pick what fits the role.
\namesection{Abdurrahman}{Shaffar}{
  \urlstyle{same}\href{https://shaffar.com}{shaffar.com} |
  \href{mailto:abdurrahman@shaffar.com}{abdurrahman@shaffar.com} |
  \href{https://scholar.google.com/citations?user=ATQK5gcAAAAJ}{Google Scholar}
}

%%%%%%%%%%%%%%%%%%%%%%%%%%%%%%%%%%%%%%
%
%     COLUMN ONE (LEFT)
%
%%%%%%%%%%%%%%%%%%%%%%%%%%%%%%%%%%%%%%
\begin{minipage}[t]{0.33\textwidth}

%%%%%%%%%%%%%%%%%%%%%%%%%%%%%%%%%%%%%%
%     SKILLS
% NOTE(hlim): List ONLY skills you are genuinely proficient in.
% ATS systems keyword-scan this section — use the exact terms from the
% job description (e.g., "ROS2", "PyTorch", "C++17", "CUDA").
% Do NOT pad with tools you barely touched; it will show in interviews.
% Skills before Education is recommended for industry applications.
% Reference: https://www.sweresume.app/articles/faang-resume-guide/
%%%%%%%%%%%%%%%%%%%%%%%%%%%%%%%%%%%%%%
\section{Skills}
\subsection{Programming}
C++ \textbullet{} Python \textbullet{} CUDA \textbullet{} CMake
\sectionsep

\subsection{Frameworks \& Tools}
Qiskit \textbullet{} QuTIP \textbullet{} Git \textbullet{} Docker
\sectionsep

\subsection{Research}
Quantum Microwave Communication \textbullet{} Continuous-Variable Quantum Information \textbullet{} Entanglement Distribution \textbullet{} Rydberg Atomic Receivers \textbullet{} Wireless Communication Systems \textbullet{} Quantum-Classical System Modeling
\sectionsep

\subsection{Languages}
Indonesian (Native) \textbullet{} English (Proficient) \textbullet{} Korean (Basic)
\sectionsep

%%%%%%%%%%%%%%%%%%%%%%%%%%%%%%%%%%%%%%
%     LINKS
% NOTE(hlim): Replace with your own links.
% Scholar is optional — include it only if you have notable citations.
%%%%%%%%%%%%%%%%%%%%%%%%%%%%%%%%%%%%%%
\section{Links}
Website:// \href{https://shaffar.com}{\bf shaffar.com} \\
GitHub:// \href{https://github.com/awshaff}{\bf awshaff} \\
LinkedIn:// \href{https://linkedin.com/in/awachid}{\bf awachid} \\
Scholar:// \href{https://scholar.google.com/citations?user=ATQK5gcAAAAJ}{\bf Abdurrahman Wachid Shaffar}
\sectionsep

%%%%%%%%%%%%%%%%%%%%%%%%%%%%%%%%%%%%%%
%     KEY OPEN-SOURCE PROJECTS
% NOTE(hlim): This section is your coding portfolio.
% Open-source projects with GitHub stars demonstrate two things:
%   (1) you can write clean, reusable, production-quality code, and
%   (2) others depend on your work — star count is a proxy for adoption.
% Keep your strongest 3-5 projects here. If you have no open-source
% work yet, start one related to the role you're targeting.
% Format: project name + star badge + one-line description + venue/year.
%%%%%%%%%%%%%%%%%%%%%%%%%%%%%%%%%%%%%%
\section{Key Open-Source Projects}
\small{Projects others rely on, reflecting both coding quality and real-world impact.}
\vspace{0.3cm}
\begin{tightemize}
\item \href{https://github.com/MIT-SPARK/KISS-Matcher}{\textbf{KISS-Matcher}} \githubstars{643} \\
Out-of-the-box fast and robust global point cloud registration for large-scale mapping. (\textit{ICRA'25})
\item \href{https://github.com/MIT-SPARK/VGGT-SLAM}{\textbf{VGGT-SLAM}} \githubstars{882} \\
Dense real-time RGB SLAM on SL(4) manifold; globally consistent 3D reconstruction. (\textit{NeurIPS'25})
\item \href{https://github.com/url-kaist/patchwork-plusplus}{\textbf{Patchwork++}} \githubstars{842} \\
$>$100\,Hz LiDAR ground segmentation; deployed in AV and mobile robot production. (\textit{IROS'22})
\item \href{https://github.com/LimHyungTae/ERASOR}{\textbf{ERASOR}} \githubstars{494} \\
${\sim}$94\% dynamic-object artifact removal, enabling long-term autonomous navigation. (\textit{RA-L'21})
\end{tightemize}
\sectionsep

%%%%%%%%%%%%%%%%%%%%%%%%%%%%%%%%%%%%%%
%     EDUCATION
% NOTE(hlim): Education goes LAST for industry applications.
% Once you have 2+ years of experience, degree details are rarely
% a deciding factor. Keep it brief: degree, school, year.
% GPA is worth including only if it is exceptional (e.g., 4.0/4.0).
%%%%%%%%%%%%%%%%%%%%%%%%%%%%%%%%%%%%%%
\section{Education}

\subsection{Kyung Hee University (KHU)}
\descript{Ph.D. in Electronics and Information Convergence Engineering}
\location{Mar. 2022 -- Feb. 2028 | Yongin, South Korea}
Advisor: Prof. Hyundong Shin \\
Dissertation: Quantum Microwave Communication
\sectionsep

\subsection{Universitas Gadjah Mada (UGM)}
\descript{B.S. in Physics}
\location{Aug, 2015 -- Nov. 2019 | Yogyakarta, Indonesia}
Advisor: Prof. Pekik Nurwantoro \\
Dissertation: Quantum Mechanics Implementation in Quantum Cryptography
\sectionsep

\end{minipage}
\hfill
%%%%%%%%%%%%%%%%%%%%%%%%%%%%%%%%%%%%%%
%
%     COLUMN TWO (RIGHT)
%
%%%%%%%%%%%%%%%%%%%%%%%%%%%%%%%%%%%%%%
\begin{minipage}[t]{0.66\textwidth}

%%%%%%%%%%%%%%%%%%%%%%%%%%%%%%%%%%%%%%
%     EXPERIENCE
% NOTE(hlim): This is the most important section.
% Every bullet must follow the XYZ formula:
%   "Accomplished [X] as measured by [Y], by doing [Z]."
%
% BAD:  "Worked on SLAM system for autonomous vehicles."
% GOOD: "Reduced localization failure rate by 40% in corridor environments
%        by implementing degeneracy-aware weighting in the odometry pipeline."
%
% Quantify everything. If you cannot measure directly, estimate scope:
%   "deployed on 3+ platforms", "city-scale", "10+ robots", "~90% reduction".
% Power words: real-time, large-scale, robust, deployed, production, pipeline.
% NOTE(hlim): Do NOT list Teaching Assistant (TA), grading, or academic service roles here.
% Industry hiring managers look for engineering and research impact in the Experience section.
% TA duties (office hours, grading, course prep) do not signal the skills tech companies hire for.
% To show mentorship or leadership instead, frame it as a bullet inside your research role:
%   e.g. "Mentored 5+ graduate researchers, contributing to 4+ published papers."
%%%%%%%%%%%%%%%%%%%%%%%%%%%%%%%%%%%%%%
\section{Experience}

\runsubsection{SPARK Lab, MIT}
\descript{| Postdoctoral Associate}
\location{Apr. 2024 -- Present | Cambridge, MA}
\vspace{0.3cm}
\begin{tightemize}
\item Reduced multi-session map inconsistency in city-scale environments by building a large-scale multi-robot SLAM system enabling real-time 3D mapping across 10+ robots and sessions.
\item Accelerated research output by mentoring 4+ graduate researchers, contributing to 4+ published papers on 3D perception and autonomous mapping.
\item Grew combined GitHub adoption to 4K+ stars by maintaining \href{https://github.com/MIT-SPARK/TEASER-plusplus}{TEASER++} \githubstars{2.1K}, \href{https://github.com/MIT-SPARK/Hydra}{Hydra} \githubstars{773}, \href{https://github.com/MIT-SPARK/KISS-Matcher}{KISS-Matcher} \githubstars{643}, \href{https://github.com/MIT-SPARK/VGGT-SLAM}{VGGT-SLAM} \githubstars{882}.
\end{tightemize}
\sectionsep

\runsubsection{Korea Advanced Institute of Science and Technology (KAIST)}
\descript{| Postdoctoral Fellow}
\location{Mar. 2023 -- Mar. 2024 | Daejeon, South Korea}
\begin{tightemize}
\item Improved localization robustness in dynamic scenes by developing real-time LiDAR-inertial and visual SLAM pipelines, deployed on 3+ mobile robot platforms.
\item Reduced dynamic-object-induced map corruption by developing deep learning perception modules for moving object segmentation in crowded and unstructured environments.
\end{tightemize}
\sectionsep

\runsubsection{StachnissLab, University of Bonn}
\descript{| Visiting Scholar}
\location{Nov. 2022 -- Feb. 2023 | Bonn, Germany}
\begin{tightemize}
\item Cut dynamic-object artifacts in LiDAR maps by ${\sim}$90\% by building ERASOR2, an instance-aware removal system for static map building in urban scenes (RSS 2023).
\end{tightemize}
\sectionsep

\runsubsection{NAVER LABS}
\descript{| Research Intern}
\location{Apr. 2022 -- Sep. 2022 | Gyeonggi-Do, Korea}
\begin{tightemize}
\item Achieved robust global LiDAR registration tolerating up to 99\% outlier correspondences by developing Quatro, exploiting ground segmentation for reliable loop closure in degenerate urban environments (ICRA 2022 / IJRR 2023).
\end{tightemize}
\sectionsep

%%%%%%%%%%%%%%%%%%%%%%%%%%%%%%%%%%%%%%
%     KEY RELEVANT PROJECTS
%
% NOTE(hlim): *** CUSTOMIZE THIS SECTION FOR EVERY APPLICATION. ***
%
% This is the core of a project-driven industry resume.
% Unlike the open-source portfolio on the left (which stays fixed),
% this section should change depending on the role:
%
%   Applying to AV perception team?
%     → Highlight 3D detection, sensor fusion, segmentation projects.
%   Applying to mapping / localization team?
%     → Highlight SLAM, HD mapping, loop closure projects.
%   Applying to robot learning / autonomy team?
%     → Highlight RL, imitation learning, sim-to-real projects.
%
% Each bullet must follow the XYZ formula (see Experience section above).
% The GitHub link proves the code exists and is public. Stars prove adoption.
%
% DUMMY ENTRIES BELOW — replace with your own relevant projects:
% Reade this article: https://www.sweresume.app/articles/faang-resume-guide/
%%%%%%%%%%%%%%%%%%%%%%%%%%%%%%%%%%%%%%
\section{Selected Projects}
\vspace{0.3cm}
\begin{tightemize}
% Replace project name, GitHub URL, star count, description, and venue.
\item \href{https://github.com/your-org/lidar-slam}{\textbf{[e.g., Real-Time LiDAR SLAM System]}} -- Reduced end-to-end localization drift by X\% in city-scale environments by implementing degeneracy-aware factor graph optimization, running at 10\,Hz on embedded hardware (\textit{[Team/Company, Year]})
\item \href{https://github.com/your-org/3d-detector}{\textbf{[e.g., 3D Object Detection Pipeline]}} -- Improved long-range detection recall by X\% by fusing LiDAR and camera with a late-fusion Transformer, processing at 20\,Hz on NVIDIA Orin (\textit{[Team/Company, Year]})
\item \href{https://github.com/your-org/hd-map}{\textbf{[e.g., Incremental HD Map Builder]}} -- Scaled occupancy mapping to campus-scale ($>$1\,km$^2$) by designing an incremental voxel compression pipeline, cutting storage by 60\% vs. naive accumulation (\textit{[Team/Company, Year]})
\end{tightemize}
\sectionsep

%%%%%%%%%%%%%%%%%%%%%%%%%%%%%%%%%%%%%%
%     PUBLICATIONS
% NOTE(hlim): For industry, publications are a credibility signal — NOT the story.
% One line is enough. Recruiters do not read paper titles.
% The venue names (ICRA, NeurIPS, RA-L) signal research quality.
% If you have 3 or fewer papers, list them individually with a one-line description.
%%%%%%%%%%%%%%%%%%%%%%%%%%%%%%%%%%%%%%
\section{Publications}
10\texttt{+} first-author papers at top-tier robotics and AI venues: IEEE ICRA, RSS, IEEE RA-L, NeurIPS, IEEE/CVF ICCV, IJRR
\sectionsep

%%%%%%%%%%%%%%%%%%%%%%%%%%%%%%%%%%%%%%
%     AWARDS
% NOTE(hlim): Keep only the most impressive 3-5 awards.
% Cut anything that needs explanation to sound impressive.
% Competition wins (especially ranked ones) are more compelling than
% certificates of participation or minor honorable mentions.
%%%%%%%%%%%%%%%%%%%%%%%%%%%%%%%%%%%%%%
\section{Awards}
\begin{tabular}{rl}
2025 & Outstanding Reviewer Award, IEEE ICRA'25 (5 awardees out of 7,641 reviewers) \\
2024 & RSS Pioneers 2024 (30 selected out of 202 candidates) \\
2024 & 1st Prize, HILTI SLAM Challenge, IEEE ICRA'24 \\
2023 & Best Paper Award, IEEE RA-L (5 of 1,100 papers) \\
2023 & 1st Prize, HILTI SLAM Challenge, IEEE ICRA'23 (63 international teams) \\
2023 & CES'23 Innovation Award (via technology transfer with HILLS Robotics) \\
2022 & 2nd Prize, HILTI SLAM Challenge, IEEE ICRA'22 \\
2020 & Student Best Paper Award, ICCAS \\
\end{tabular}
\sectionsep

\end{minipage}
\end{document}
